\section{Limitations}

First, EHA currently uses a synthetic mini-web. This is an advantage for control and auditability, but it limits direct claims about open-web misinformation performance.

Second, the current main run contains 100 tasks. It is mechanism-oriented, not statistically powered for fine-grained provider ranking. Wilson intervals in \Cref{tab:main-model-decomposition} help show sampling uncertainty, but they do not convert the run into a provider-ranking study.

Third, the benchmark has a fixed schema and scoring contract. Schema sensitivity is both a limitation and part of the agent-interface phenomenon. The generated-lore clarified-schema mini-rerun shows that field wording can change role assignment, but this does not retroactively replace the frozen main-run score.

Fourth, active-verification scoring still needs more human audit. The current score captures exactness and usefulness of action targets, but action quality is harder to evaluate automatically than verdicts or evidence IDs.

Fifth, this paper does not claim that the evaluated models will behave the same way on real institutions, real legal or medical tasks, or live open-web misinformation. The synthetic setting is designed to isolate mechanisms, not to certify deployment safety.

Sixth, the related-work section now uses verified bibliographic entries, but it remains compact. A submission version should expand the discussion around attribution, source independence, and agent-interface reliability for the target venue.
